%!TEX root = ../dissertation.tex
% !TeX spellcheck = en_GB

\chapter{Conclusion}
\label{conclusion}

Starting with Chapter \ref{chapt:2}, this thesis presented and studied current state-of-the-art algorithms proposed and used in the field of black-box optimization. The path follows the historical progress and evolution of the field showing ideas and techniques used by authors to improve the fundamental idea of Nesterov \cite{nesterov2017random}, to enhance performances in optimization of non-convex functions and Reinforcement Learning problems. The main basic ideas are the following:
\begin{itemize}
	\item to parallelize the computations made by the algorithm with the use of multi-core CPUs;
	\item to specialize the choice of random directions in order to take into account the exploration-exploitation trade-off, view taken from the Reinforcement Learning framework;
	\item to implement a novel way to estimate the gradient and adaptivity of the parameters of the algorithm itself.
\end{itemize}

In Chapter \ref{chapt:3} a novel algorithm was presented inspired by the ones presented in the literature. Its adaptivity and use of previous gradients information allow to improve the performance with respect to state-of-the-art methods. In ASHGF the adaptivity of the parameters and the novel way of estimating the gradient, was adapted to the handling of the trade-off between exploration and exploitation. 

Chapter \ref{chapt:4} shows numerical results of ASHGF and other state-of-the-art algorithms on a widely used testbed and two Reinforcement Learning problems. The presented algorithm outperforms the others in all benchmarks: in terms of data and performance profiles, it achieves a better accuracy most of the time while requiring less function evaluations; in terms of its average behaviour, it is the one that performs better on average, outperforming all the algorithms in some functions; solving Reinforcement Learning problems it achieves higher average rewards faster than the competitors.

Further work should focus on rewriting the algorithm, and the others, in order to parallelize its execution and then testing them in a greater pool of functions and, more importantly, reinforcement learning problems, which in this case were few due to technical limitations.